\begin{question}{25}{
    Het Space Operations Centre (SpOC) van de NAVO bewaakt een constellatie van militaire communicatiesatellieten (MILSATCOM) die cruciale data-uplinks verzorgen voor grondtroepen in een operatiegebied. 
    Vijandelijke actoren passen regelmatig RF-jamming (radiofrequentie-storing) toe om deze uplinks tijdelijk te onderbreken.

    Uit operationele data over de afgelopen 18 maanden blijkt dat vijandelijke jamming-aanvallen gemiddeld $\lambda=3.2$ aanvallen per dag voorkomen volgens een Poissonproces.
    Elke aanval duurt gemiddeld $25$ minuten met een standaardafwijking van $3.5$ minuten, en verstoort de communicatieverbinding volledig voor de duur van de aanval. 

    Het SpOC beschikt over een Counter-Space Response Team (CSRT) dat actief tegenmaatregelen kan inzetten (frequentiehoppen, uplink-switching).
    Het team kan per dag maximaal $5$ aanvallen effectief neutraliseren. 
    Aanvallen boven deze capaciteit leiden tot een communicatieblackout voor de grondtroepen.
}
    \subquestion{5}{
        Bereken de kans dat het SpOC in een willekeurige week \emph{minimaal} $20$ jamming-aanvallen registreert.
    }
    \solution{
        
    }

    \subquestion{5}{
        Bereken de kans dat het CSRT op een willekeurige dag overbelast raakt. 
    }
    \solution{
        
    }

    \subquestion{5}{
        De NAVO-norm stelt dat de kans op een communicatieblackout per etmaal kleiner dan \SI{5}{\percent} moet zijn. 
        Bepaal de minimale CSRT-capaciteit (in aantallen aanvallen per dag die geneutraliseerd kunnen worden) waarmee aan deze norm wordt voldaan.
    }
    \solution{

    }

    \subquestion{}

\end{question}